\noindent
\begin{tabular}{ll}
\textbf{Thesis}          & Interpretable Embodied AI for a Hybrid Bi-pedal/Quadrupedal Robot \\
                         & Using Sparse Autoencoders and Vision-Language Model Feature Steering \\
\textbf{Student}         & Mr.~Tanakorn Youngmeesuk \\
\textbf{Student ID.}     & 65011563 \\
\textbf{Degree}          & Bachelor of Engineering \\
\textbf{Program}         & Robotics and AI Engineering \\
\textbf{Year}            & 2026 \\
\textbf{Thesis Advisor}  & Prof.~Dr.~Poom Konghuayrob \\
\end{tabular}

\vspace{1cm}

\begin{center}
\textbf{\Large ABSTRACT IN ENGLISH}
\end{center}

\vspace{0.5cm}

This thesis presents an end-to-end embodied artificial intelligence pipeline for a
hybrid bi-pedal/quadrupedal robot whose locomotion mode can be switched on demand.
The central challenge addressed is twofold: enabling a robot with multiple locomotion
configurations to accept and execute natural-language commands without requiring any
programming knowledge from the operator, and making the internal decision-making
process of the underlying Vision-Language Model transparent and auditable.

The AI brain is built upon Qwen3-VL-2B-Thinking, a compact two-billion-parameter
multimodal model fine-tuned on a purpose-built dataset of 10{,}000 robot-conversation
records generated through an automated multi-threaded pipeline. Fine-tuning teaches
the model to produce structured \texttt{call\_robot()} tool-call outputs from RGB
scene images and free-form user requests. To expose the model's internal
representations, a TopK Sparse Autoencoder (SAE) is trained on 100{,}000 layer-20
hidden-state activations, decomposing the model's 2{,}048-dimensional residual stream
into 32{,}768 interpretable, near-mono\-semantic feature dimensions. Feature-level
steering is demonstrated by injecting SAE decoder column vectors directly at the
identified layer, modifying robot behaviour at inference time without any model
retraining. Experiments confirm that Feature~\#5086 in the agentic SAE model acts as
a primary task-planning neuron: steering this feature reliably elicits structured
coordinate-sequence outputs across a range of pick-and-place, stacking, and
drawer-manipulation tasks. The complete system is deployed through a Gradio web
interface with a real-time steering-strength slider, making interpretable robot
control accessible to non-technical users.

\vspace{0.4cm}
\noindent\textbf{Keywords:} Vision-Language Model, Sparse Autoencoder, Mechanistic
Interpretability, Feature Steering, Hybrid Robot, Embodied AI, Natural-Language
Robot Control